\documentclass[12pt]{article}
\usepackage[utf8]{inputenc}
\usepackage[T1]{fontenc}
\usepackage{hyperref}
\title{Sinasina Language Resources for K-12 Curriculum Development}
\author{Compiled by Claude Sonnet 5.5 (claude-sonnet-5-5) for Andr\'as Kornai}
\date{}
\begin{document}
\maketitle

\section{Introduction}
Sinasina (ISO 639-3 sst; Glottocode sina1271) is a Trans-New Guinea language of Simbu Province, Papua New Guinea. The name is an outsider label for several Chimbu--Wahgi varieties in Tabare Rural LLG. Speakers use local names such as Tabare, Kere and Kondo. The 2000 census counted about 21,000 native speakers (other sources give higher figures). Documentation found in search is thin.

\section{Primary Educational Resources}
The Wikipedia article reports that literacy primers were produced, credited to Rosalie Ranquist and Nell Dreghorn. I could not verify this, and I found no copies: \url{https://en.wikipedia.org/wiki/Sinasina_language}.

\section{Linguistic and Reference Materials}
\begin{itemize}
\item Wikipedia, Sinasina language (the claim of a dictionary compiled by Charles Turner is unsourced and unconfirmed): \url{https://en.wikipedia.org/wiki/Sinasina_language}
\item ASJP word list page for Sinasina: \url{https://asjp.clld.org/languages/SINASINA_2}
\end{itemize}

\section{Religious and Community Resources}
The Wikipedia article reports a New Testament in the Tabare dialect printed in 1975 (New Tribes Mission, Charles Turner). It is unverified here. Joshua Project has a people-group profile: \url{https://joshuaproject.net/index.php/people_groups/14190}.

\section{Digital and Audio Resources}
The University of Hawai'i ScholarSpace holds an item, apparently a Tabare word-list recording from the Kaipuleohone archive. I did not check its contents: \url{https://scholarspace.manoa.hawaii.edu/items/9b69c510-eddd-443e-aa9d-4501b5b5ce16}.

\section{Curriculum Development Context}
No school use of Sinasina was found. A curriculum would have to be built on limited primers and a possible dictionary and New Testament that need locating through library and SIL or Wycliffe catalogues. Dialect variation (Tabare, Kere, Kondo) must be addressed.

\end{document}